% 模板作者：Zydragon-x
% 联系作者：5509121090@email.ncu.edu.cn

\documentclass{beamer}
\usepackage{hyperref}
\usepackage{xeCJK}
\setCJKmainfont{AR PL UKai TW}
% 這兩行一定要加，中文才能自動換行
\XeTeXlinebreaklocale "zh" 
\XeTeXlinebreakskip = 0pt plus 1pt



\usepackage[T1]{fontenc}
\usepackage{style}

% other packages
\usepackage{latexsym,amsmath,xcolor,multicol,booktabs,calligra}
\usepackage{graphicx,pstricks,listings,stackengine}
\usepackage{mathrsfs}
\usepackage{mathtools} 
\usepackage{extarrows} 
\usepackage{amsthm}
\usepackage{CJKutf8}
\usepackage{amssymb}
\usepackage{amsmath}
\usepackage{import}
\usepackage{xifthen}
\usepackage{pdfpages}
\usepackage{import}
\usepackage{transparent}
\usepackage{overpic}
\usepackage{setspace}
\usepackage{rotating}
\usepackage{caption}
\usepackage{wrapfig}
\usepackage{svg}
\usepackage{dsfont}
\usepackage{subcaption}
\usepackage{xcolor}

%\usetheme{Copenhagen}

%\newcommand{\sixedge}{\mbox{\parbox[t][-2pt][c]{10pt}{\scriptsize$\blacktriangledown\hspace{-3.6pt}\blacktriangle\hspace{-4pt}\blacktriangledown$}\hspace{-10.2pt}\parbox[b][10.7pt][c]{10pt}{\scriptsize$\blacktriangle\hspace{-3.6pt}\blacktriangledown\hspace{-3.6pt}\blacktriangle$}}}

\author{郝嘉誠、徐樂融、洪偉傑、蘇士紘、蘇彥維 \\ 指導老師：林偉傑(台大數學)、李志煌(台大數學)}
\title{Planar Statistical Physics and \\Bernoulli Percolation}
\subtitle{Group 7 Midterm Report}
\institute{NCTS USRP}
\date{Jul 22, 2022}

\newcommand{\incfig}[1]{%
\import{./}{#1.pdf_tex}
}
\newcommand{\sixedge}{\mbox{\parbox[t][-2pt][c]{10pt}{\scriptsize$\blacktriangledown\hspace{-3.6pt}\blacktriangle\hspace{-4pt}\blacktriangledown$}\hspace{-10.2pt}\parbox[b][10.7pt][c]{10pt}{\scriptsize$\blacktriangle\hspace{-3.6pt}\blacktriangledown\hspace{-3.6pt}\blacktriangle$}}}

% defs
\def\cmd#1{\texttt{\color{red}\footnotesize $\backslash$#1}}
\def\env#1{\texttt{\color{blue}\footnotesize #1}}
\definecolor{deepblue}{rgb}{0,0,0.5}
\definecolor{deepred}{rgb}{0.6,0,0}
\definecolor{deepgreen}{rgb}{0,0.5,0}
\definecolor{halfgray}{gray}{0.55}

\lstset{
    basicstyle=\ttfamily\small,
    keywordstyle=\bfseries\color{deepblue},
    emphstyle=\ttfamily\color{deepred},    % Custom highlighting style
    stringstyle=\color{deepgreen},
    numbers=left,
    numberstyle=\small\color{halfgray},
    rulesepcolor=\color{red!20!green!20!blue!20},
    frame=shadowbox,
}




\begin{document}



\begin{frame}
    
    \titlepage

\end{frame}

%\begin{frame}
    %\tableofcontents[sectionstyle=show,subsectionstyle=show/shaded/hide,subsubsectionstyle=show/shaded/hide]
%\end{frame}

\begin{frame}{Table of Content}
    \tableofcontents
\end{frame}

\section{Settings}
%\begin{frame}{Test}
%    \begin{block}{Text}
%    hihi
%    \end{block}
%\end{frame}
\begin{frame}{Motivation}
    \begin{exampleblock}{} % 加 * 
        \begin{itemize}
            \item For connection or flowing problem, we can observe cellular network, polymer growth, or simply drip coffee.
            \item It will be interesting to know how the probability of an edge being open, $p$, affects phenomena in macro scale. To model, we may use a simpler mesh to start.
            \item For $p = 0$ (fully closed) and $p = 1$ (fully open) case, the result is obvious, but what happen for $p$ in between?
        \end{itemize}
        \begin{figure}[htbp]
            \centering
            \includegraphics[width=0.3\linewidth]{drip_coffee.jpg}
            \captionsetup{labelformat=empty}
            \caption{Drip coffee }
            
        \end{figure}
    \end{exampleblock}
    \
\end{frame}

\begin{frame}{Sample Space}
    \begin{exampleblock}{} % 加 * 
        \begin{itemize}
            \item The sample space $\Omega$ is $\{0,1\}^E$, containing all the possible configurations of edges being open/closed.\\
            \item A \textbf{configuration} $\omega \in \Omega$ gives a connecting situation on edges. "1" means open, and "0" means closed. \\
            \item A \textbf{cluster} means a connected component in $\omega$, and infinite cluster is a cluster of infinite 'size'. 
    \end{itemize}
    \end{exampleblock}
    \begin{figure}[htpb]
        \centering
        \includegraphics[width=0.23\linewidth]{p1.png}
        \captionsetup{labelformat=empty}
        \caption{A square lattice.}
    \end{figure}
\end{frame}

\begin{frame}{Bernoulli Percolation Model}
    \begin{exampleblock}{} % 加 * 
        In Bernoulli percolation model, each edge is open with probability $p$, and closed with $1-p$. The edges are sampled independently.
        %We consider the probaility space $(\{0,1\}^E, \mathcal{F},\mathbb{P}_p)$
        \begin{itemize}
            %\item $\mathcal{F}$ is the $\sigma$-algebra generated by events depending on finitely many edges (cylinder set).
            \item $\mathbb{P}_p$ is the probability measure on events, edges in which are open with probability $p$. 
    \end{itemize}
    \end{exampleblock}
    \begin{figure}[htpb]
        \centering
        \includegraphics[width=0.2\linewidth]{p0.55.png}
        \captionsetup{labelformat=empty}
        \caption{A configuration sampled at $p = 0.55$}
    \end{figure}
\end{frame}

%\begin{frame}{Definition}
%    \begin{itemize}
%        \item $\mathbb{P}_p \colon$measures the probability of $\omega_p\in\{0,1\}^E$ when each edge is open with probability $p$.
%        \item If a vertex is in an infinite cluster, then it can be "connected to $\infty$."
%    \end{itemize}
%\end{frame}

%\begin{frame}{Percolation Function}
%    \begin{itemize}
%        \item Define 
%        \begin{align*}
%        \theta \colon [0,1] & \longrightarrow \mathbb{R}\\
%        p & \longmapsto \mathbb{P}_p(\text{0 belongs to an infinite cluster}).
%        \end{align*}
%        \item If $\theta(p)=0$, then then $0$ cannot be connected to $\infty$ almost surely.
%        \item Question:\\
%        When is $\theta(p)>0$?
%    \end{itemize}
%\end{frame}

\begin{frame}{Double Phase Transition}
    \begin{itemize}
        \item Consider the graph $\mathbb{T}_d\square\mathbb{Z}$. 
        \item Questions: $p_c$? Uniqueness of infinity clutser?
        \item $0<p_c<\textcolor{red}{$p_u$}<1$
    \end{itemize}
\end{frame}

%\begin{frame}{Critical Probability}
%    \begin{itemize}
%        \item If $p\in [0,1]$ is greater than the "critical probability", then there exists an infinite cluster with positive probability.
%        \item $p_c$ is defined as 
%        \[
%        \sup\{p\in[0,1] \mid \theta(p)=0\}.
%        \]
%    \end{itemize}
%    \begin{figure}
%            \centering
%            \subcaptionbox[some text]{ A sampling for $p=0.45$.   }[.45 \textwidth]
%            {\includegraphics[width = .45\linewidth]{p=0.45.png}
%            }
%            \hfill
%             \subcaptionbox[some text]{ A sampling for $p=0.55$.   }[.45 \textwidth]
%            {\includegraphics[width = .45\linewidth]{p=0.55.png}
%            }                
%        \end{figure}
%\end{frame}

% \begin{frame}{Coupling}
%     \begin{itemize}
%         \item How to compare the ``size" of $\omega_p$ and $\omega_{p'}$ with $p\leq p'$? 
%         \item Let $\{\mu_e\}_{e \in E(G)}$ be i.i.d Unif$(0,1)$ random variables.
%         \item For each edge $e$, we define
%         \[\omega_p(e) = \begin{cases}
%             1 &, \mu_e \leq p; \\
%             0 &, \mu_e > p.
%         \end{cases}\]
%         \item If $\omega_p(e)=1$, then $\omega_{p'}(e)=1$.
%         \item Intuitively, $\omega_{p'}$ has more open edges.
%     \end{itemize}
% \end{frame}

\section{Self-Avoiding Walk}

%洪偉傑
\begin{frame}{Self-Avoiding Walk}
    \begin{center}
        \Large Self-Avoiding Walk
    \end{center}
    \begin{enumerate}
        \item Definition
        \item Existence of the connective constant for a graphs
        \item Connective constant of some good graphs
        \item Conclusion
    \end{enumerate}
\end{frame}
\begin{frame}{Self avoiding walk}

Definition : For a graph $G,$
\begin{figure}[h]
    \centering
    \includesvg[width=8cm]{SAW1.svg}
\end{figure}
\end{frame}
\begin{frame}{Existence of the connective constant for a graph}
Let $\lambda_n=\#$ SAWs of length $n$ from $O$ in $G.$
    Observe:
    \[
    \lambda_{n+m}\leq \lambda_n\cdot \lambda_m.
    \]
\begin{figure}[h]
    \centering
    \includesvg[width=8cm]{SAW2.svg}
\end{figure}
\end{frame}
\begin{frame}{Existence of the connective constant for a graphs}
    We use the fact that $\lambda_{n+m}\leq \lambda_n\cdot\lambda_m,$ we can deduce that 
    \[
    \mu=\mu(G):=\lim_{n\to\infty}\sqrt[n]{\lambda_n}.
    \]
    exists.\\[5pt]
    The constant $\mu$ is called \textbf{connective constant.}
\end{frame}
\begin{frame}{Connective constant of some good graphs}
    \begin{figure}[htp]
        \centering
        \includesvg[width=10cm]{conclusion.svg}
    \end{figure}
\end{frame}
\begin{frame}{Connective constant of some good graphs}
    1. Tree ($\mathbb{T}_3$)
    \begin{wrapfigure}{l}{5.2cm}
    \flushleft
    \includesvg[width=5cm]{T3.svg}
\end{wrapfigure}
    $\lambda_n=\lambda_{n-1}\times 2\Rightarrow \mu(\mathbb{T}_3)=2.$\\
    \vspace{5cm}
    Similarly, $\mu(\mathbb{T}_d)=d-1,\ \forall d\geq 2.$
\end{frame}
\begin{frame}{Connective constant of some good graphs}
    2. hexagonal mesh $\sixedge$
    \begin{wrapfigure}{l}{5.2cm}
    \flushleft
    \includesvg[width=5cm]{hexagonalmesh.svg}
\end{wrapfigure}
    $\lambda_{n+1}\leq \lambda_n\times 2\Rightarrow \mu(\sixedge)\leq 2.$\\
    Moreover, $\lambda_{n+5}\leq \lambda_n\times 30\Rightarrow\mu(\sixedge)\leq \sqrt[5]{30}.$\\[2cm]
    Actually, $\mu(\sixedge)=\sqrt{2+\sqrt{2}}.$ (H. Duminil-Copin, S. Smirnov, 2010)
\end{frame}
\begin{frame}{Connective constant of some good graphs}
    3. Triangular mesh $\triangle$
    \begin{wrapfigure}{l}{5.2cm}
    \flushleft
    \includesvg[width=5cm]{nothing.svg}
\end{wrapfigure}
Let $S_n$ be \# SAWs of length $n$ that only move to the following four directions in every steps. We get
\[
\lambda_{n}\geq S_n=3S_{n-1}+2S_{n-2}.
\]
Thus,
\[
\mu(\triangle) \geq \lim_{n\to\infty} \sqrt[n]{S_n}=\frac{3+\sqrt{13}}{2}.
\]
\end{frame}
\begin{frame}{Connective constant of some good graphs}
    4. $\mathbb{Z}\times\{0,1\}$
    \begin{wrapfigure}{l}{5.2cm}
    \flushleft
    \includesvg[width=5cm]{ladder.svg}
\end{wrapfigure}
If we consider $S_n=\#$ SAWs of length $n$ that only move $\uparrow,\rightarrow,\downarrow$ in every step, we deduce that 
\[
\lambda_n\geq S_n=S_{n-1}+S_{n-2},\ n\geq 2.
\]
Thus, 
\[
\mu(\mathbb{Z}\times\{0,1\})\geq\lim_{n\to\infty}\sqrt[n]{S_n}=\frac{1+\sqrt{5}}{2}.
\]
Actually, $\mu(\mathbb{Z}\times\{0,1\})=\cfrac{1+\sqrt{5}}{2}\ .$
\end{frame}
\begin{frame}{Conclusion}
\begin{center}
        It is hard to compute $\mu.$
\end{center}
\end{frame}
\begin{frame}{Conclusion}
    \[
    \mu(\mathbb{Z}^2)\leq 4,\quad \mu(\mathbb{Z}^3)\leq 6,\cdots  \mu(\mathbb{Z}^d)\leq 2d.
    \]
    Recall that 
    \begin{align*}
         &p_c=\sup\{p\, | \, \theta(p)=0\}\\
         \Rightarrow &p_c(\mathbb{Z}^d)\geq\frac{1}{\mu(\mathbb{Z}^d)}\geq \frac{1}{2d}>0.
    \end{align*}
\end{frame}
\begin{frame}{Bond percolation and Site percolation}
    \begin{figure}[htp]
        \centering
        \includesvg[width=\textwidth]{boundandsitepercolation.svg}
    \end{figure}
\end{frame}
\begin{frame}{Critical Probability}
    Given $G$ be a graph, recall that 
    \begin{align*}
        &\theta(p)=\mathbb{P}_p[O\Leftrightarrow\infty\mbox{ in }G]\\
         &p_c=\sup\{p\, | \, \theta(p)=0\}
    \end{align*}
    We study $p_c$ for some nice graph.
\end{frame}
\begin{frame}{Critical Probability of some nice graphs}
    1. Tree ($\mathbb{T}_3$)
    \begin{figure}[htp]
        \centering
        \includesvg[width=6.5cm]{T3.svg}
    \end{figure}
\end{frame}

% 蘇士紘
% 證明 p_c 在平面網格上會大於等於二分之一
\section{Some Results on $\mathbb{Z}^2$}
\begin{frame}{$p_c(\mathbb{Z}^2)\geq\frac{1}{2}$}
    \begin{block}{Theorem. (Kesten)}
    $p_c(\mathbb{Z}^2)=\frac{1}{2}$
    \end{block}
    % \begin{enumerate}
    %     \item $0 < p_{c} < 1$ on $\mathbb{Z}^{d}$
    %     \item Dual graph on $\mathbb{Z}^{2}$
    %     \item $\theta(p_{c}) = 0$ on $\mathbb{Z}^{2}$
    %     \item $p_{c} = \frac{1}{2}$ on $\mathbb{Z}^{2}$
    % \end{enumerate}
    First show that $\theta(\frac{1}{2})=0$. (Hence $p_c(\mathbb{Z}^{2})\geq \frac{1}{2}$)\\
    Method: Dual graph argument $+$ The square root trick. 
    \begin{figure}[htp]
    \centering
    \includesvg[width=10cm]{Ex3.svg}
    \end{figure}
\end{frame}


%\begin{frame}{Dual graph on $\mathbb{Z}^{2}$}
%    \begin{center}
%    \includegraphics[width=5cm]{./dool1.png}
%    
%    An example of a graph $G$ and its dual graph $G^{*}$.
%    \end{center}
%\end{frame}

\begin{frame}{Dual graph on $\mathbb{Z}^{2}$}
    \begin{center}
    \includegraphics[width=4cm]{./dool3.png}
    
     Overlapping a graph $G$ and its dual graph $G^{*}$.
    \end{center}
\end{frame}

%\begin{frame}{Harris-FKG Inequality}
%\begin{block}{Definition (Partial order on $\Omega$, $\leq$)}
%For $\omega,\omega^\prime\in\Omega$, $\omega\leq\omega^\prime\Leftrightarrow\omega(e)\leq\omega^\prime(e)$ $\forall e\in %E(G)$.
%\end{block}
%\begin{block}{Definition (Increasing events)}
%An event $A$ is said to be increasing if for all $\omega,\omega'\in\Omega$,
%$$\omega\leq\omega^\prime\;and\;\omega\in A\Rightarrow\omega^\prime\in A.$$
%\end{block}
%\end{frame}

\begin{frame}{Harris-FKG Inequality}
\begin{block}{Theorem (Harris-FKG inequality)}
Let $A,B$ be two increasing events. Then $\mathbb{P}(A\cap B)\geq\mathbb{P}(A)\mathbb{P}(B)$.
\end{block}
$Intuition.$ $\mathbb{P}(A \mid B)\geq\mathbb{P}(A)$ for increasing events $A$ and $B$.
\begin{block}{Corollary (The square-root trick)}
Let $\{A_i\}_{i=1}^n$ be a family of increasing events. Then 
\[\max_{1 \leq i \leq
 n} \mathbb{P}(A_i) \geq 1 - \left[1 - \mathbb{P}\left(\bigcup_{i = 1}^n A_i\right)\right]^{\frac{1}{n}}.\]
\end{block}
\end{frame}

%\begin{frame}{$\theta(P_{c}) = 0$ on $\mathbb{Z}^{2}$}
%    reference where $d \geq 11$ is proved
%    
%    Mean-field behavior for nearest-neighbor percolation in $d \geq 11$.
%    Robert Fitzner, Remco van der Hofstad.    
%\end{frame}

\section{Tools for Bernoulli Percolation}

% 蘇彥維
% \begin{frame}{Tools}
%     \begin{enumerate}
%         \item Harris - FKG inequality
%         \item van den Berg - Kesten inequality
%         \item Russo's Identity
%         % \item 
%     \end{enumerate}
% \end{frame}

\begin{frame}{p_c(\mathbb{Z}^2)\leq\frac{1}{2}}
    \begin{block}{Theorem (Exponential decay property).}
    Let $G=\mathbb{Z}^d$, $d\geq 2$.\\
    If $p<p_c(G)$, then there exists $c(p)>0$ s.t. for all $n\geq 1$,$$
    \mathbb{P}_p(0 \xlongleftrightarrow{\text{$\Lambda_n$}} \partial\Lambda_n)\leq e^{-c(p)n}.
    $$
    If $p>p_c(G)$, then $\mathbb{P}_p(0\longleftrightarrow\infty)\geq\frac{p-p_c}{p(1-p_c)}$.
    \end{block}
    $Application.$ $p_c(\mathbb{Z}^2)\leq\frac{1}{2}$.
    
    % \textit{Sketch of proof of }$p_c(\mathbb{Z}^2)\leq\frac{1}{2}$.\\
    % Assume $\frac{1}{2}<p_c$.\\
    % Dual graph observation $+$ bound using exponential decay.\\
    % Get contradiction: hence $p_c\leq\frac{1}{2}$.
\end{frame}

\begin{frame}{Russo-Seymour-Welsh theory}
    \begin{block}{Theorem}
    Let $\rho>0$. There exists $c=c(\rho)>0$ s.t. for all $n\geq 1$, we have
    $$
    c\leq\mathbb{P}_{\frac{1}{2}}\big[\mathcal{H}(\rho n, n)\big]\leq 1-c,
    $$
    where $\mathcal{H}(\rho n, n)$ denotes the horizontal crossing event in a rectangle with size $[0,\rho n]\times [0,n]$.
    \end{block}
    \begin{itemize}
        \item Once we get a $c(\rho)$ for some $\rho>1$, we can obtain $c(\rho')$ for any $\rho'>1$.
        \item Example:
        \begin{equation*}
            \mathbb{P}_\frac{1}{2}\big[\mathcal{H}(4n,n)\big]\geq\mathbb{P}\big[\mathcal{H}(2n,n)\big]^5\geq c(2)^5.
        \end{equation*}
        \item Example (Combining $\mathcal{H}(n,n)$):
        \begin{equation*}
            \mathbb{P}_\frac{1}{2}\big[\mathcal{H}(3n,2n)\big]\geq\frac{1}{128}.
        \end{equation*}
    \end{itemize}
\end{frame}

\begin{frame}{Russo-Seymour-Welsh theory}
    % \begin{block}{Example}
    Consider a simply connected domain with a smooth boundary $\Omega$ with distinct boundary points $a,b,c,d$.\\
    \begin{itemize}
        \item Generate a finite graph $\Omega^\delta=\delta\mathbb{Z}^2\cap\Omega$ for $\delta>0$
        \item Define $a^\delta, b^\delta, c^\delta, d^\delta\in\Omega^\delta$ to be the closest points to $a,b,c,d\in\partial\Omega$
        \item Define $(a^\delta b^\delta),\,(c^\delta d^\delta)$ be the paths on $\partial\Omega^\delta$ from $a^\delta$ to $b^\delta$, from $c^\delta$ to $d^\delta$ counterclockwise.
    \end{itemize}
    % \end{block}
    
    \begin{figure}[b]
        \includegraphics[width=5cm]{omega.png}
        \includegraphics[width=5cm]{omega_2.png}
    \end{figure}

\end{frame}



% \begin{frame}{van den Berg-Kesten inequality}
% \begin{block}{Definition}
% (Witness) Given an event $A$ and $\omega\in A$. The edge set $I(\omega)\subseteq E(G)$ is a $witness$ of $A$ for $\omega$ if\\
% 1. $\omega=1$ on $I(\omega)$.
% \\
% 2. $\omega'=\omega$ on $I(\omega)\Rightarrow \omega'\in A$.\\
% (Realized disjointly) Given $\omega\in\Omega$. Two events $A$ and $B$ can be \textit{realized disjointly} if there exists witnesses for $A$ and $B$, $I(\omega)$ and $J(\omega)$, such that $I(\omega)\cap J(\omega)=\emptyset$. We collect such $\omega$ for $A$ and $B$ as $A\circ B$.
% \end{block}
% \begin{figure}[htbp]
%             \centering
%             \includegraphics[width=0.66\linewidth]{fig_realised_disjointly.png}
%             \captionsetup{labelformat=empty}
%             % \caption{Drip coffee \cite{wiki:perco}}
            
% \end{figure}
% \end{frame}

% \begin{frame}{van den Berg-Kesten inequality}
% \begin{block}{Theorem (BK inequality)}
% For two increasing events $A,B$ depending only on finitely many edges, we have $\mathbb{P}(A\circ B)\leq\mathbb{P}(A)\mathbb{P}(B)$.
% \end{block}
% $Intuition.$ Negative correlation. If the event $B$ can not use the edges used by $A$, then the probability for $B$ to hold gets smaller.
% \begin{figure}[htbp]
%             \centering
%             \includegraphics[width=0.66\linewidth]{fig_realised_disjointly.png}
%             \captionsetup{labelformat=empty}
%             % \caption{Drip coffee \cite{wiki:perco}}
            
% \end{figure}
% \end{frame}

% \begin{frame}{Russo's identity}
% \begin{block}{Definition (pivotal edge)}
% For $\omega\in\Omega$ and an event $A$, $e\in E$ is a pivotal edge for $A$ and $\omega$ if
% $$
% \omega(e)=1\Rightarrow\omega\in A\text{ and }\omega(e)=0\Rightarrow\omega\notin A.
% $$
% \end{block}
% \begin{block}{Theorem (Russo's Identity)}
% If $A$ is an increasing event depending on finitely many edges, then $$\frac{d}{dp}\mathbb{P}_p(A)=\mathbb{E}_p(N(A))$$,
% where $N(A)$ is the number of the pivotal edges for the event $A$.
% \end{block}

%     %\begin{figure}[htbp]
%     %    \centering
%     %    \includegraphics[width=0.56\linewidth]{sausages.png}
%     %    \captionsetup{labelformat=empty}
%         % \caption{Drip coffee \cite{wiki:perco}}
%     %\end{figure}
    
% \end{frame}



\section{Prospect}

\begin{frame}{Future Directions}
    \begin{exampleblock}{1) On a more general percolation model (e.g. non i.i.d.)}
    \begin{itemize}
        \item Do we have RSW property, exponential decay or other inequalities on the random-cluster model?
        % \item Site Bernoulli percolation on triangle lattice ($p_c = \frac{1}{2}$)
    \end{itemize}
    \end{exampleblock}
    \begin{exampleblock}{2) Bernoulli percolation on other graphs}
        \begin{itemize}
            \item Periodic lattices (triangle, hexagon etc.)
            % \item High dimensional lattice
            \item Cayley graph (fintely generated groups)\\
                Relationship between percolation properties with group properties.
        \end{itemize}
    \end{exampleblock}
    
    
\end{frame}

%\section{Reference}

%\begin{frame}[allowframebreaks]
%    \bibliography{ref}
%    \bibliographystyle{alpha}
    % 如果参考文献太多的话，可以像下面这样调整字体：
    % \tiny\bibliographystyle{alpha}
%\end{frame}

%\begin{frame}
%    \begin{center}
%        {\Huge\calligra Thanks!}
%    \end{center}
%\end{frame}

%\begin{frame}{Frame Title}
    
%\end{frame}

\end{document}